\documentclass[10pt]{beamer}

% Tema e pacotes essenciais
\usetheme{Madrid}
\usecolortheme{default}
\usefonttheme{professionalfonts}

\usepackage[utf8]{inputenc}
\usepackage[T1]{fontenc}
\usepackage[brazil]{babel}
\usepackage{amsmath,amssymb,amsthm,amsfonts}
\usepackage{bm}

% Comandos usuais
\newcommand{\Prb}{\mathbb{P}}
\newcommand{\E}{\mathbb{E}}
\newcommand{\Rif}{\mathrm{Rif}}
\newcommand{\TV}{\mathrm{TV}}

% Título e metadados
\title[Riffle Shuffle]{Mistura de cartas e o modelo \emph{Riffle Shuffle}}
\author{Jair Donadelli}
\institute{Universidade Federal}
\date{}

\begin{document}

%==================================================
\begin{frame}
  \titlepage
\end{frame}

%==================================================
\begin{frame}{Modelo e distribuição do \emph{riffle shuffle}}
  Seja \(S_n\) o grupo de permutações de \(\{1,\dots,n\}\). 
  No modelo de \emph{riffle shuffle invertido} (\emph{inverse riffle}), 
  atribuímos a cada carta um rótulo \(b_i\in\{0,1\}\), i.i.d. com probabilidade \(1/2\),
  e então ordenamos o baralho por ordem não-decrescente dos rótulos, 
  preservando a ordem relativa das cartas com o mesmo rótulo (ordenamento estável).

  Assim, para cada \(\sigma\in S_n\),
  \[
    \Rif(\sigma) 
    = \frac{1}{2^n}\#\bigl\{ b\in\{0,1\}^n : 
      \text{a ordenação estável de }b\text{ produz }\sigma \bigr\}.
  \]
  Essa é a medida de probabilidade de um \emph{riffle shuffle}.

  \vspace{0.5em}
  O suporte da medida \(\Rif\) é o conjunto \(R\subset S_n\) das permutações
  obtidas intercalando duas subsequências crescentes.
  O número de tais permutações é
  \[
    |R| = 2^n - n.
  \]
\end{frame}

%==================================================
\begin{frame}{Contagem dos \emph{riffle shuffles}}
  \textbf{Proposição.} 
  O número de permutações em \(S_n\) que podem ser obtidas
  como interleavings de duas subsequências crescentes é \(2^n - n\).

  \vspace{0.5em}
  \textbf{Esboço de prova.} 
  Cada vetor binário \(b\in\{0,1\}^n\) define uma permutação \(f(b)\)
  ao intercalar, de modo estável, as cartas com rótulo \(0\) e as com rótulo \(1\).
  Duas sequências binárias distintas produzem a mesma permutação
  se e somente se todos os zeros aparecem antes de todos os uns,
  isto é, quando \(b=0^k1^{n-k}\) para algum \(k\in\{0,\dots,n\}\).
  Há \(n+1\) dessas sequências redundantes,
  e portanto
  \[
    |R| = 2^n - (n+1) + 1 = 2^n - n.
  \]
  \(\square\)
\end{frame}

%==================================================
\begin{frame}{Distribuição de um passo}
  Para \(n\ge1\), temos
  \[
    \Rif(\sigma)
    = \frac{1}{2^n}\#\{b\in\{0,1\}^n : \text{ordenação estável}(b)=\sigma\}.
  \]

  Em particular:
  \[
    \Rif(\mathrm{id}) = \frac{n+1}{2^n},
  \]
  pois a identidade é obtida pelas \(n+1\) sequências \(0^k1^{n-k}\),
  \(k=0,\dots,n\).
  
  \vspace{0.5em}
  As demais permutações \(\sigma\in R\setminus\{\mathrm{id}\}\)
  têm probabilidade \(\Rif(\sigma) = m_\sigma/2^n\),
  onde \(m_\sigma\ge1\) é o número de rótulos binários que as geram.
\end{frame}

%==================================================
\begin{frame}{Tempo de mistura e lema do tempo de parada forte}
  Seja \((X_k)_{k\ge0}\) a cadeia de Markov gerada por \(\Rif\),
  e \(U\) a distribuição uniforme sobre \(S_n\).
  Se \(T\) é um tempo de parada forte tal que,
  condicionalmente a \(T=j\),
  \(X_j\) é uniforme e independente de \(T\),
  então
  \[
    \|Q^{(k)}-U\|_{\TV} \le \Prb[T>k].
  \]
  (Prova: argumento clássico do lema do tempo de parada forte.)
\end{frame}

%==================================================
\begin{frame}{Distribuição do tempo de parada}
  Em cada embaralhamento, cada carta recebe um rótulo uniforme
  em \(\{0,1,\dots,2^k-1\}\).
  O evento \(\{T\le k\}\) ocorre quando todos os \(n\) rótulos são distintos,
  de modo que
  \[
    \Prb[T>k]
      = 1 - \prod_{i=1}^{n-1}\left(1-\frac{i}{2^k}\right).
  \]
  Assim,
  \[
    \|Q^{(k)} - U\|_{\TV} 
       \le 1 - \prod_{i=1}^{n-1}\left(1-\frac{i}{2^k}\right).
  \]
\end{frame}

%==================================================
\begin{frame}{Aproximação pelo paradoxo do aniversário}
  Para \(m\ll 2^{k/2}\),
  \[
    \prod_{i=1}^{m-1}\left(1-\frac{i}{2^k}\right)
      \approx \exp\!\left(-\frac{m(m-1)}{2^{k+1}}\right).
  \]
  Assim, o tempo \(k\) necessário para que 
  \(\|Q^{(k)} - U\|_{\TV}\) seja pequeno
  satisfaz aproximadamente
  \[
    2^k \approx n^2,
    \quad\text{ou seja,}\quad
    k \approx \tfrac{3}{2}\log_2 n.
  \]
  Para \(n=52\), isso corresponde a cerca de \(7\) embaralhamentos.
\end{frame}

%==================================================
\begin{frame}{Referência clássica}
  \begin{thebibliography}{9}
    \bibitem{BayerDiaconis}
      D.~Bayer e P.~Diaconis,
      \emph{Trailing the Dovetail Shuffle to Its Lair},
      Annals of Applied Probability, vol.~2, no.~2 (1992), pp.~294–313.
  \end{thebibliography}
\end{frame}

\end{document}
